\section{Experimental Methods for Archived Builds}
\label{sec:supp-methods}

\noindent\textbf{Build lineage and deployment.}
Results retain their original build identifiers (Table~\ref{tab:supp-builds}). The
\emph{continuation and accounting baseline}, committed as \texttt{d99ec92},
implements the layered continuation and the accounting transitions; its
compensation commits together with its representation. The artifact records
per-binary SHA-256 fingerprints and the configuration of each run. The
nine-service continuation, compensation and repayment, storage, and
short-throughput series use this build. The initial fourteen-service paired
runs use the same single-obligation rule before the batch-result tag was
separated from the direct-result tag, and the cross-type codec regressions
cover that separation. Members in the baseline storage, in-memory
continuation, and short-throughput controls run with GOGC=200 and
GOMAXPROCS=16; this override is not applied to committee, gateway, or wallet
processes.

The \emph{decision-first build} (accounting domain
\texttt{DIRECT\_ACCOUNTING\_V5\_DECISION\_FIRST}) implements resubmission from
public certificates, the compensation decision, and the separate
representation commands. It derives from the branch base \texttt{0e7bae5}, a
documentation-only descendant of \texttt{d99ec92}. The nine functional runs
use the initial implementation snapshot, whose source-package SHA-256 begins
\texttt{2c91b0251024322c}. The later snapshot of this series, whose package SHA-256 begins
\texttt{449b4d4072c99588}, rejects a decision whose position does not match
committed state before reading any historical block and raises the bounded
due-decision scan from 8 to 32 entries per round; the ordinary payment path is
unchanged. The six ordinary-path runs, the full test suite, and the race tests
use the later snapshot of this series. The audited commit \texttt{f185f73} additionally
requires the certificates attached to a payment to correspond one-to-one to
its certificate inputs, and it does not change the payments accepted in the
measured workloads. Numbers from the baseline and decision-first builds are
reported separately.

\begin{table}[t]
  \caption{Builds and the series they support.}
  \label{tab:supp-builds}
  \centering
  \small
  \begin{tabular}{@{}p{0.2\columnwidth}p{0.3\columnwidth}p{0.4\columnwidth}@{}}
    \toprule
    Build & Identifier & Series \\
    \midrule
    Continuation and accounting baseline
      & \texttt{d99ec92}
      & 100-hop chains, synchronous-member and storage controls,
        compensation and repayment matrix, paired interference runs,
        short throughput. \\
    Decision-first
      & Base \texttt{0e7bae5}; functional package \texttt{2c91b0251024322c};
        later series snapshot \texttt{449b4d4072c99588}; audited \texttt{f185f73}
      & Nine functional runs, six ordinary-path runs, regression and race
        tests. \\
    Archived evaluation
      & \texttt{721800c}
      & 100-hop chains, the complete N/R/C/P matrix, and historical-reading controls (Section~S2). \\
    Partial-approval reclamation
      & \texttt{bf71c4d}
      & Conflict and restart regressions; 1,000-payment smoke and 100-hop continuation. \\
    Archived
      & \texttt{e3870e1}/\texttt{a083ff6}, \texttt{f0c0a274},
        \texttt{8e1fb40}/\texttt{5807546}, \texttt{bb65af69}/\texttt{13be856},
        \texttt{f82f8cd}, \texttt{c1ebf2d}
      & Operating range, original continuation, capital, organization
        paths, batch integration, source-failure checkpoint. \\
    \bottomrule
  \end{tabular}
\end{table}

The archived operating-range runs record base \texttt{e3870e1} and
implementation \texttt{a083ff6}. Original continuation records
\texttt{f0c0a274}, capital records base \texttt{8e1fb40} and implementation
\texttt{5807546}, and organization paths record base \texttt{bb65af69} and
verified source \texttt{13be856}. The archived batch integration uses
\texttt{f82f8cd}, and the structural repair counts were repeated on the
baseline code. The source-failure checkpoint \texttt{c1ebf2d} uses the
superseded late-source rule and does not support the current claims.

Mac trials share one M4 Max host with 16 CPU cores and 64\,GiB RAM, and the
structural microbenchmarks run on Windows amd64. Baseline memory chains and
short throughput trials use in-memory service state and NoSync wallets.
Compensation and repayment trials use disk-backed committee application state
and BlockStore with in-memory members and gateway. Storage controls change
only member commits between bbolt NoSync and synchronous transactions. Each
run starts from fresh genesis. Six additional fourteen-service disk-backed
chain runs use the baseline accounting rule and user-paid FUEL. Their three
fast 100-hop runs take 668.960--688.080\,ms, and the inter-hop confirmation
runs take 114.496--118.742\,s. They form a separate storage and deployment
configuration from the nine-service results in the main paper.

\noindent\textbf{Decision-first functional runs.}
The nine runs use the fourteen-service configuration with 450 initial
outputs, a disk-backed committee, members in their default NoSync mode, flush
and gossip settings of 10\,ms, and explicitly authorized organization-funded FUEL; signatures and all amount
and consumption checks remain enabled. Ordinary-payment stress wallets use the
NoSync wallet option, and the two-hop fault demonstrations use the default
wallet mode. In each run, an injected fault prevents the original publisher
from submitting the parent payment or running INSTALL for it after its full QC exists, so
the original signers hold their approvals but not the QC. The child spends the
parent's TXCer. In the resubmission treatment the driver never submits the
original payment. The control treatment disables the automatic resubmission
step with an experiment switch. Ordinary backup submission by members
holding installed complete material remains enabled; INSTALL itself
replicates that material within the organization. The
paused treatment also disables resubmission and, while a marker file exists,
makes the input and part adaptation HTTP endpoints of all four committee
members return a paused status; consensus, decisions, and payment processing
are not paused.

Times start at the first HTTP send of the child payment. The fast receipt ends
when the wallet has verified the certificate and recorded the output. Public
events are observed by polling the followers' status endpoints, so they
include polling delay, and the difference between two public observations is
not a protocol delay measured from the exact commit. The compensation time
includes the configured 30-s obligation period, the advance of public time by
committee ticks, and background scheduling. In the paused treatment the script
submits the late source after the compensation and waits until the status
endpoints of all four committee members show the obligation repaid and the
representation neither committed nor materialized. It stores these four
observations, removes the marker, and then waits for the representation to be
committed and materialized at all four members. The status endpoint's
stability flag compares the header hash, data hash, and transaction identity.
The materialization step and the real-consensus integration tests check the
complete block identity and part-set header, so the flag alone is not offered
as the complete identity proof. The repayment and the resumption of adaptation
have no timestamps of their own. The measurements contain no comparison with
periodic scanning, vote collection, or the earlier coupled build under paused
adaptation.

Each treatment has three rounds, and every functional run passes four checks:
equal state hashes at the four committee members, a zero gap, closed
payments, and drained outboxes. The resubmission and control treatments each
carry three additional runs of 300 ordinary payments. Per-run values are in
Table~\ref{tab:supp-functional}.

\noindent\textbf{Decision-first ordinary-path runs.}
Six runs compare the build at the branch base with the final decision-first
snapshot in the order AB, BA, AB. Each sends 300 payments at 100 payments/s
with 64 concurrent fast requests and at most 256 outstanding, using NoSync
stress wallets. The elapsed time includes the public confirmation of all 300
payments and member following; it is neither a single-payment time nor a
maximum throughput. The send span is about 2.990\,s in every run. Summary
values in the main text are medians of the per-run statistics, not quantiles
pooled over payments. Ordinary final outputs carry no input certificates, so
these runs do not measure the storage retained for payments whose inputs are
not final. Per-run values are in Table~\ref{tab:supp-cost}.

\begin{table*}[t]
  \caption{Per-run records of the nine decision-first functional runs.
    Times are client observations from the child payment's wallet send;
    ``Parent fast'' is the original payment's own receipt. Public columns
    include polling. A dash marks a column that does not apply or whose
    time was not recorded. ``Auto.\ source'' does not apply in the control
    and paused rows because automatic resubmission is disabled there; the
    driver later submits the source, which executes and repays the issuer
    (verified by the repayment audit), but that instant was not
    timestamped. In the resubmission rows no compensation or
    materialization occurs.}
  \label{tab:supp-functional}
  \centering
  \small
  \setlength{\tabcolsep}{4pt}
  \begin{tabular}{@{}llrrrrrr@{}}
    \toprule
    Treatment & Run & \shortstack{Child fast\\(ms)} & \shortstack{Parent fast\\(ms)}
      & \shortstack{Child public\\(s)} & \shortstack{Auto. source\\(s)} & \shortstack{Compensation\\(s)}
      & \shortstack{Materialized\\(s)} \\
    \midrule
    Resubmission & 1 & 13.94 & 9.02 & 1.414 & 2.014 & -- & -- \\
     & 2 & 13.94 & 8.94 & 1.990 & 2.565 & -- & -- \\
     & 3 & 13.91 & 9.02 & 1.440 & 2.039 & -- & -- \\
    \addlinespace
    Control & 1 & 11.96 & 8.97 & 1.364 & -- & 35.839 & 36.994 \\
     & 2 & 12.94 & 8.97 & 1.413 & -- & 35.888 & 37.021 \\
     & 3 & 12.01 & 9.00 & 1.388 & -- & 35.864 & 37.019 \\
    \addlinespace
    Paused & 1 & 12.00 & 7.98 & 1.362 & -- & 35.837 & 42.748 \\
     & 2 & 13.92 & 9.04 & 1.341 & -- & 35.814 & 42.736 \\
     & 3 & 13.92 & 9.03 & 1.841 & -- & 35.891 & 42.798 \\
    \bottomrule
  \end{tabular}
\end{table*}

\begin{table*}[t]
  \caption{Per-run records of the six ordinary-path runs. Run orders are
    AB, BA, AB (A $=$ branch base, B $=$ decision-first).
    Quantiles are over the 300 payments of one run.}
  \label{tab:supp-cost}
  \centering
  \small
  \setlength{\tabcolsep}{4pt}
  \begin{tabular}{@{}llrrrrrr@{}}
    \toprule
    Build & Run & \shortstack{Elapsed\\(s)} & \shortstack{Fast P50\\(ms)} & \shortstack{Fast P95\\(ms)}
      & \shortstack{Block P50\\(ms)} & \shortstack{Dispatch P95\\(ms)} & \shortstack{Send span\\(s)} \\
    \midrule
    Branch base & 1 & 3.871 & 1.112 & 27.441 & 1151.0 & 0.338 & 2.990 \\
     & 2 & 3.815 & 1.096 & 17.951 & 1101.1 & 0.295 & 2.990 \\
     & 3 & 3.805 & 1.052 & 39.945 & 1076.0 & 0.437 & 2.990 \\
    \addlinespace
    Decision-first & 1 & 4.000 & 1.071 & 30.530 & 1127.9 & 0.430 & 2.990 \\
     & 2 & 3.845 & 1.073 & 22.084 & 1103.9 & 0.355 & 2.990 \\
     & 3 & 3.791 & 1.045 & 33.744 & 1052.3 & 0.383 & 2.990 \\
    \bottomrule
  \end{tabular}
\end{table*}

\noindent\textbf{Regression coverage.}
Compensation decisions are rejected when they are not yet due, when the
network, slot, or amount does not match, when the reserve is insufficient, or
when they conflict with the recorded decision. A repeated decision has no
economic effect, including after repayment, where it returns no changes,
effects, or fee outputs, and both orders, source first and decision first, are
covered. Resubmission is tested for cross-organization and multi-level
sources, repeated exposure, and payments already observed or already queued,
which are not rebuilt. Ordinary verification rejects mismatched material, and
the original approval, the payment bytes, and the CAL quota are unchanged
before and after resubmission. Representation requires a successful decision,
remains possible after the source's repayment, and does not restore member
progress or Worker slices a second time. For identity and replay, real
threshold adaptation preserves the complete consensus identity, the successor
body, and pre-signed successor references. Replay of the original history
reproduces the application state at every height, and item-wise and
differently grouped representation of one decided set give the same final
body and part openings, with 1, 2, 4, and 8 inputs spanning two real parts.
The full test suite and static checks pass at the later snapshot of this series on Windows
and macOS, the race detector passes on the member, representation, and
committee-runtime packages on macOS, and the full suite and static checks pass
again at the audited commit.

\noindent\textbf{Timing and aggregation.}
Certified-receipt latency generally runs from the first wallet HTTP send
through recipient verification and local atomic recording. Capital and
source-failure drivers include a small amount of pre-send construction, so
their comparisons stay within the same driver family. A local NoSync
transaction completion is not a synchronous durability receipt. Public and
member completion timestamps record observation of authenticated execution,
not the committee's internal commit. The sender's single monotonic clock
measures cross-site ACK latency, including the return path, and a recipient
can start its successor before the preceding sender receives that ACK.
Offered targets, actual send rates, and completion rates including drain are
retained separately. Captions distinguish transaction quantiles from
statistics across runs, and within-run sample variability is not an estimate
of cross-run uncertainty.

\noindent\textbf{Delivery-gate runs.}
On the archived delivery-gate build, the gated variant waits for three
members to INSTALL complete material or for authenticated public success,
whichever occurs first, and collects no second signature quorum. At 200
payments/s for 30\,s, three paired runs per mode each send and complete 6,000
payments. The median of run-level receipt P50s is 1.300\,ms with immediate
delivery and 2.167\,ms with the gate. Separate 100-hop trials reduce median
chain time from 287.983 to 200.356\,ms (Fig.~\ref{fig:delivery-gate}).

\begin{figure*}[!t]
  \centering
  \includegraphics[width=\textwidth]{delivery-gate.pdf}
  \caption{Delivery-gate effect (ms). Left: run-level receipt P50 at
    200 payments/s, 6,000 payments/run, three paired repetitions.
    Right: whole-chain time for three 100-hop chains per mode. Columns
    show the median of the three run values; dots show individual runs.}
  \label{fig:delivery-gate}
\end{figure*}

\noindent\textbf{Baseline-build controls.}
The reserve trajectories of the three 5\% runs of the compensation and
repayment matrix (Fig.~\ref{fig:reserve-recovery}) are sampled at the
driver's release schedule, so the lowest sampled reserve is not a
reserve-sizing result.

\begin{figure}[t]
  \centering
  \includegraphics[width=\columnwidth]{recovery-reserve.pdf}
  \caption{Sampled reserve trajectories for the three 5\% runs (60
    compensated sources each). Repayment restores the initial 60,000 CAL.
    The vertical axis is restricted to 59,600--60,000 CAL to show temporary
    use. Samples reflect the driver's release schedule.}
  \label{fig:reserve-recovery}
\end{figure}

Three pairs of fourteen-service runs, ordered normal/withheld,
withheld/normal, and normal/withheld, use the initial baseline build, before
the batch-result tag separation, with committee and member databases on disk
and member and wallet NoSync. Each sends 4,000 payments at 100 payments/s for
40\,s. Alongside them, four additional two-hop demonstrations have sources
that are delivered normally (normal run) or withheld, compensated after
30\,s, and later repaid (withheld run). All 4,000 payments complete in all six
runs. The final CAL held by each owner role (30,000 and 830,000) and both
reserve balances ($10^{12}$) coincide across the six runs, and the withheld
runs record 400\,CAL paid and 400\,CAL repaid, so the final allocation is
consistent with delay neutrality. Receipt P50s for normal and withheld runs
are 1.070/1.127, 1.089/1.108, and 1.083/1.097\,ms, and P95s are
26.84/59.23, 29.49/58.44, and 27.41/55.15\,ms; completion rates are
97.2--97.8 payments/s in all runs. Receipt P50 is 1.3--5.3\% higher and P95 is
1.98--2.21 times the paired normal values. Concurrent compensation therefore
affects the receipt tail in this configuration, and these pairs do not
isolate the responsible stage.

Six runs of 1,000 payments at 100 payments/s on the nine-service
configuration (committee on disk, wallets NoSync) differ only in the member
store. With NoSync members, receipt P50 is 1.002--1.035\,ms and P95
33.97--40.62\,ms, with 93.11--93.93 payments/s completed including drain.
With synchronous member commits, P50 is 43.19--51.85\,ms and P95
80.98--96.56\,ms, with 90.99--92.72 payments/s. All payments complete. Because
wallets remain NoSync, this isolates the cost of synchronous member commits.

\noindent\textbf{Representation microbenchmarks and batch integration.}
These structural measurements use the earlier coupled representation command,
which also carried the economic step. Microbenchmarks on Windows amd64 with
Go~1.27.1 place $k\in\{1,2,4,8\}$ missing inputs in one historical block and
part. Across five repeats, single commands require $k$ commands and $k$ part
adaptations, whereas a batch uses one of each, and input adaptations remain
$k$. At $k=8$, total encoded command bytes fall from 127,336 to 2,752
(Fig.~\ref{fig:repair-cost}). The dispersed control separates batching from
compact encoding: eight independent target blocks still require eight
commands and eight part adaptations, while bytes fall from 30,008 to 5,440. At
$k=1$, both formats perform one part adaptation and encode 3,751 versus 680
bytes. Compact fields therefore explain part of the byte reduction, and
shared block representation explains the reduced adaptation count. These are
encoded command sizes, not total network traffic or total cryptographic work.
They also exclude the 88-byte decision command that each decided output
carries in the decision-first build.

\begin{figure*}[!t]
  \centering
  \includegraphics[width=\textwidth]{repair-cost.pdf}
  \caption{Same-block, same-part representation costs of the earlier
    coupled command versus obligation count $k$. Panels show
    command/part-adaptation counts and total encoded command payloads in
    KiB (log scale). Five repeats per configuration give identical
    structural measurements. Input adaptations remain $k$.}
  \label{fig:repair-cost}
\end{figure*}

An archived fourteen-service, disk-backed batch-integration run
(\texttt{f82f8cd}) combines four actual missing sources into a single
1,568-byte batch whose BatchID agrees across all committee replicas. After
late-parent handling, all 1,000 subsequent payments offered at 100 payments/s
complete. Receipt P50/P95 are 1.048/26.531\,ms, and total completion time is
11.120\,s, including observation and shutdown. Gaps and outboxes drain, and
committee state hashes agree. This batch carried economic steps and predates
the separation of decisions from representation.

\noindent\textbf{Externally funded authorization controller.}
The capital experiments sample member availability every 200\,ms. Let $\rho$
be the declared CAL authorization demand per second, counting both parent and
child requests in each workload unit, and let $u$ count unique requests whose
certification remains unfinished. The controller sets
\begin{align*}
 b &= 200+100u,\\
 H_{\mathrm{low}} &= \rho\,(0.8\,\mathrm{s})+b,\\
 H_{\mathrm{target}} &= 1.25H_{\mathrm{low}}.
\end{align*}
Here the watermarks and burst allowance are in CAL. When the minimum observed
member availability $A_{\min}$ falls below $H_{\mathrm{low}}$, the proposed
grant increment is
\[
 \Delta B=100\left\lceil
 \frac{1.5(H_{\mathrm{target}}-A_{\min})}{100}
 \right\rceil,
\]
subject to the configured authorization cap and available external funding. At
most one increment is in flight. Its public command debits the external source
and credits the reserve and grant atomically, and following members add the
resulting share difference without clearing existing approvals. The experiment
supplies $\rho$ in advance rather than estimating unknown production demand.
Controller decisions are requests for actual funded transfers, not a mechanism
for creating backing capital.

\noindent\textbf{Capital workload accounting.}
Each condition has one 300-s run with one Worker per member and sufficient
user FUEL. Starting from 3,600 CAL, the constant, rate-step, and
delayed-parent conditions end at 14,400, 21,600, and 23,300 CAL after external
transfers of 10,800, 18,000, and 19,700 CAL
(Table~\ref{tab:eval-capital}). Each completes 6,000 two-hop units, or 12,000
payments, spending 1,128,000 user FUEL. Corresponding fixed-high controls
complete, while the fixed-low condition closes 46 of 6,000 units, leaves 5,826
unstarted, and records 64 partially approved payments. These categories have
different units and are not summed as one payment count. The per-run tables
preserve the full seven-condition comparison without implying three
repetitions per condition. The series ran on an earlier build in which no
compensation occurred.

\begin{table}[t]
  \caption{Actual capital additions. Each condition has one 300-s run,
    starts with 3,600 CAL, and completes 12,000 payments. Both columns
    are CAL; every addition debits an external funded account.}
  \label{tab:eval-capital}
  \centering
  \small
  \begin{tabular}{@{}lrr@{}}
    \toprule
    Workload & Final reserve & Added capital \\
    \midrule
    Constant: 20 units/s & 14,400 & 10,800 \\
    Step: 10 to 30 units/s & 21,600 & 18,000 \\
    Parent delayed by 3\,s & 23,300 & 19,700 \\
    \bottomrule
  \end{tabular}
\end{table}

\noindent\textbf{Identities, resource use, and retained history.}
The multi-identity configurations use shared clients, so 2,048 wallet
identities do not represent 2,048 separate operating system processes.
Comparing independent inputs with ten-hop payments, the three-run medians are
12.67 versus 15.24 core-ms per payment and 59.32 versus 71.83\,KiB of HTTP
payload per payment. Memory reporting separates total node RSS and retained
history from active obligations, pending work, and outboxes, so an increasing
historical footprint is not reported as an increasing payment backlog.
Per-run resource summaries and their archived source locations are included in
the artifact index.

\noindent\textbf{Endpoint-specific Lightning reference.}
The archived deployment uses LND v0.21.3-beta and Bitcoin Core 31.1 in
regtest, with synchronous LND persistence and its 50-ms channel-commit
setting. Each of three fresh deployments has two nodes and a direct channel.
After 20 warm-up transfers, each run issues 100 serial, alternating
10,000-satoshi payments using prepared invoices, with receiver
\texttt{SETTLED} as the completion event; the reported P50/P95 are
282.12/320.17\,ms. A separate Next run sends 300 payments at 100 payments/s
for 3\,s using in-memory service state and NoSync wallet storage.
Recipient-certified receipt has P50/P95 of 0.9625/1.2075\,ms. Workload
concurrency, durability, and completion semantics differ, so these records are
deployment references rather than a normalized acceleration ratio.

\begin{table}[t]
  \caption{Compensation decision record and representation commands (C3).
    Both are public commands committed once; they differ in what the
    stored history shows.}
  \label{tab:c3-compare}
  \centering
  \small
  \begin{tabular}{@{}p{0.2\columnwidth}p{0.37\columnwidth}p{0.37\columnwidth}@{}}
    \toprule
    & Decision record & Decision plus representation (C3) \\
    \midrule
    Economic closure
      & Reserve debit, obligation state, and fee effects commit in one
        public decision, independent of adaptation; successors unchanged.
      & The same decision; representation commands add no economic
        effect. \\
    Stored historical body
      & Original body unchanged; the funding source of a consumed input is
        read from the decision index.
      & Stored body names the reserve debit; transaction identity, owner
        authorization, and complete block identity $(H_b,P_b)$ are
        unchanged and verified. \\
    Inspecting a designated historical input
      & Consult the decision index together with the block.
      & Read the represented source under the original commitments and consult the permanent decision and current obligation state in the same $V_H$. \\
    Readers
      & An honest local replica that has verified and executed prefix $H$.
      & The same local replica, reading a committed logical revision in $V_H$; remote untrusted readers need separate authorization evidence. \\
    Storage and replay
      & Original blocks plus decision records; replay runs the original
        commands, including decisions.
      & Additionally revised bodies, revisions, and tasks; replay runs the
        same original commands and the representation commands;
        same-block commands share part adaptation and installation
        advances to the latest committed revision. \\
    Progress dependency
      & Public execution of the decision.
      & Three adaptation contributions and a current revision; economic
        closure does not depend on them. \\
    Status
      & Implemented as the economic path; measured in the decision-first
        runs.
      & Implemented; functional runs and regressions above. \\
    \bottomrule
  \end{tabular}
\end{table}

\subsection{Archived Operating-Range Series}
The archived operating-range series complements the continuation results. One
active organization uses sponsored fees, reused addresses, and independent
final inputs. Five runs of 150,000 payments each all finish
(Fig.~\ref{fig:operating}). At a target of 2,000 payments/s, two runs achieve
completion rates of 1,948.66--1,950.45 payments/s, including drain, with
receipt P50s of 2.521--2.566\,ms. Raising the target to 2,200 yields
2,103.63--2,106.29 completed payments/s. At 2,400, the achieved rate is
2,174.68, P50/P95 rise to 7.174/89.38\,ms, and the 4,096-entry pending
admission limit records 1,635 hits. The operating curve pairs achieved
completion rates with their corresponding queue pressure.

\begin{figure}[ht]
  \centering
  \includegraphics[width=0.62\columnwidth]{operating-point.pdf}
  \caption{Operating points versus target offered load. Top: completion
    throughput (payments/s), including drain. Bottom: receipt latency (ms).
    Each point is one 150,000-payment run:
    two at target 2,000, two at 2,200, and one at 2,400.
    P50/P95 summarize transactions within each run, not cross-run
    uncertainty. All services share one host.}
  \label{fig:operating}
\end{figure}
